\documentclass[fontsize=11pt, parskip=half]{scrartcl}
\usepackage[paperwidth=6.5in, paperheight=9.5in,
            margin=1cm]{geometry}
\pagestyle{plain}

\usepackage{amsmath, amssymb, amsthm, amsfonts}

\usepackage[dvipsnames]{xcolor}

\usepackage[colorlinks=true, linkcolor=blue, urlcolor=blue, citecolor=blue]{hyperref}
\usepackage{url}
\usepackage[capitalise,noabbrev]{cleveref}

\newtheorem{theorem}{Theorem}[section]
\newtheorem{lemma}[theorem]{Lemma}
\newtheorem{fact}[theorem]{Fact}
\theoremstyle{definition}
\newtheorem{definition}[theorem]{Definition}
\newtheorem{remark}[theorem]{Remark}

\crefname{section}{Problem}{Problems}
\Crefname{section}{Problem}{Problems}
\crefname{theorem}{Theorem}{Theorems}
\Crefname{theorem}{Theorem}{Theorems}
\crefname{definition}{Definition}{Definitions}
\Crefname{definition}{Definition}{Definitions}
\crefname{lemma}{Lemma}{Lemmas}
\Crefname{lemma}{Lemma}{Lemmas}
\crefname{fact}{Fact}{Facts}
\Crefname{fact}{Fact}{Facts}

\newif\ifsolutions
\solutionstrue

\usepackage{tcolorbox}
\tcbuselibrary{breakable}

\ifsolutions
  \newenvironment{solution}{%
    \begin{tcolorbox}[colback=black!5!white, colframe=black!50!white,
      title={\textbf{Solution}}, fonttitle=\normalfont, breakable]
  }{\end{tcolorbox}}
\else
  \usepackage{verbatim}
  \newenvironment{solution}{\expandafter\comment}{\expandafter\endcomment}
\fi

\tcolorboxenvironment{definition}{colback=blue!3!white, colframe=blue!40!white,
  breakable, top=4pt, bottom=4pt}
\tcolorboxenvironment{theorem}{colback=red!3!white, colframe=red!40!white,
  breakable, top=4pt, bottom=4pt}
\tcolorboxenvironment{lemma}{colback=purple!3!white, colframe=purple!40!white,
  breakable, top=4pt, bottom=4pt}
\tcolorboxenvironment{fact}{colback=orange!3!white, colframe=orange!40!white,
  breakable, top=4pt, bottom=4pt}

\newtcolorbox{defbox}[1][]{colback=blue!3!white, colframe=blue!40!white,
  title={\textbf{#1}}, fonttitle=\normalfont, breakable, top=4pt, bottom=4pt}

\newcommand{\mytitle}[1]{{\noindent\Large\textbf{#1}}}
\renewcommand{\thesection}{\arabic{section}}
\newcommand{\Msol}{\mathcal{M}_{sol}}
\newcommand{\KL}{D_{\text{KL}}}
\newcommand{\FSQ}{\mathcal{F}_{\text{SQ}}}
\newcommand{\FKL}{\mathcal{F}_{\text{KL}}}

%===================================
\begin{document}

\mytitle{Solomonoff Worksheet --- Pareto Optimality \hfill \today}

\medskip
\noindent
This is a self-contained companion to the main Solomonoff worksheet. It
establishes that the Bayesian mixture $\xi$ is \emph{uniquely Pareto-optimal}
under both KL and squared prediction loss --- no other predictor can do at
least as well as $\xi$ on every environment in the class without coinciding
with $\xi$ entirely.

Both proofs share one mechanism: a ``Pythagorean'' decomposition of the
$w$-weighted loss around $\xi$, from which uniqueness drops out as a
non-negativity statement. We treat the KL case first (\cref{prob:kl}) ---
it lives at the joint-distribution level, where the Pythagorean identity is
essentially the definition of $\xi$. The squared-loss case (\cref{prob:sq})
lives at the conditional level and needs one extra Bayes-rule step to bridge
per-history posterior weights with the prior-weighted Pareto aggregation. As
a bonus, the same identities show $\xi$ is actually the \emph{minimizer} of
the $w$-weighted loss across all predictors --- stronger than Pareto
optimality alone.

Difficulty ratings follow Knuth's scheme: \textbf{[05]} a few minutes,
\textbf{[10]} 15-minute pencil-and-paper, \textbf{[15]} half an hour or so.

%===================================
\section*{Recap of setup}

We use the same setup as the main worksheet: countable model class
$\Msol = \{\nu_1, \nu_2, \dots\}$ of environments (each $\nu$ assigning a
proper distribution $\nu(\cdot \mid x_{<t}) \in \Delta\mathbb{B}$ to every
history), prior $w_\nu > 0$ with
$\sum_\nu w_\nu = 1$, and Bayesian mixture
\[
  \xi(x) ~:=~ \sum_{\nu \in \Msol} w_\nu \, \nu(x),
  \qquad
  \xi(x_t \mid x_{<t}) ~:=~ \sum_{\nu \in \Msol} w(\nu \mid x_{<t}) \, \nu(x_t \mid x_{<t}).
\]
A \textbf{predictor} is any function $\rho$ that assigns to each history
$x_{<t}$ a distribution $\rho(\cdot \mid x_{<t}) \in \Delta\mathbb{B}$; by
chain rule, $\rho(x_{1:n}) := \prod_t \rho(x_t \mid x_{<t})$. The mixture
$\xi$ is one such predictor, and any $\nu \in \Msol$ is another.

\begin{definition}[Pareto domination]\label{def:pareto}
For a loss $\mathcal F(\nu, \rho)$ that takes an environment $\nu$ and a
predictor $\rho$, a predictor $\rho$ \textbf{weakly Pareto-dominates} $\xi$ if
\[
  \mathcal F(\nu, \rho) ~\leq~ \mathcal F(\nu, \xi)
  \qquad \text{for every } \nu \in \Msol.
\]
$\xi$ is \textbf{Pareto-optimal} (w.r.t.\ $\mathcal F$) if the only predictor
that weakly dominates it is $\xi$ itself.\footnote{This is stronger than the
usual ``no strict domination'' formulation: weak domination already forces
$\rho = \xi$, so a fortiori there is no $\rho$ that is strictly better on some
$\nu$ without being worse on another.}

Note that ``$\mathcal F(\nu, \rho) \leq \mathcal F(\nu, \xi)$ for every $\nu$''
quantifies over the class --- there is no single ``true environment.'' Each
$\nu$ takes its turn as the candidate environment; Pareto-optimality says no
$\rho$ beats $\xi$ uniformly across the whole class.

\medskip
Throughout, we fix a horizon $n \geq 1$. The two losses we study are:
\begin{align*}
  \FKL^n(\nu, \rho)
    ~&:=~
    \KL\bigl(\nu(X_{1:n}) \,\big\|\, \rho(X_{1:n})\bigr)
    ~=~
    \sum_{x_{1:n} \in \mathbb{B}^n}
      \nu(x_{1:n}) \, \ln \frac{\nu(x_{1:n})}{\rho(x_{1:n})}, \\[4pt]
  \FSQ^n(\nu, \rho)
    ~&:=~
    \sum_{t=1}^{n}
    \sum_{x_{<t} \in \mathbb{B}^{t-1}}
      \nu(x_{<t})
    \sum_{x_t \in \mathbb{B}}
      \bigl(\nu(x_t \mid x_{<t}) - \rho(x_t \mid x_{<t})\bigr)^{2}.
\end{align*}
\textbf{Histories are always summed over all of $\mathbb{B}^{t-1}$}, weighted
by the probability that the candidate environment $\nu$ reaches them: this
weight is $\nu(x_{1:n})$ for KL (baked into the joint distribution) and the
explicit factor $\nu(x_{<t})$ for the per-step squared error. So in both
cases, $\mathcal F(\nu, \rho)$ is a $\nu$-expected loss: the average prediction
error of $\rho$ when the data is generated by $\nu$. Statements below hold for
every $n$.
\end{definition}

%===================================
\section{Pareto optimality under KL loss}\label{prob:kl}

KL is the easier of the two cases. It lives entirely at the joint distribution
level $\nu(x_{1:n})$, where the Bayesian mixing identity $\sum_\nu w_\nu \,
\nu(x_{1:n}) = \xi(x_{1:n})$ is literally the definition of $\xi$. The
Pythagorean decomposition and the Pareto aggregation use exactly the same
weights, and the proof is a one-step calculation.

\bigskip

\noindent
\textbf{[10]} (Pythagorean identity for KL.) For \emph{any} predictor $\rho$
with $\rho(x_{1:n}) > 0$ on every $x_{1:n} \in \mathbb{B}^n$, show that
\begin{align*}
  &\KL\bigl(\xi(X_{1:n}) \,\big\|\, \rho(X_{1:n})\bigr)
  ~+~
  \sum_{\nu \in \Msol} w_\nu \, \KL\bigl(\nu(X_{1:n}) \,\big\|\, \xi(X_{1:n})\bigr) \\
  &=~
  \sum_{\nu \in \Msol} w_\nu \, \KL\bigl(\nu(X_{1:n}) \,\big\|\, \rho(X_{1:n})\bigr).
\end{align*}
This is an equality of three real numbers; there is no history variable
floating free --- all references to histories $x_{1:n}$ are summed over inside
the KL's, and weighted by the relevant joint distribution.

\textit{Hint.} Subtract the second sum from the third, simplify the
log-difference $\ln\tfrac{\nu}{\rho} - \ln\tfrac{\nu}{\xi} = \ln\tfrac{\xi}{\rho}$
inside the $\nu$-expectation, then swap the order of summation and use the
prior-mixing identity $\sum_\nu w_\nu \nu(x_{1:n}) = \xi(x_{1:n})$ (which holds
by definition of $\xi$ as the prior mixture on joints).

\begin{solution}
Take the difference $\sum_\nu w_\nu \KL(\nu \,\|\, \rho) - \sum_\nu w_\nu
\KL(\nu \,\|\, \xi)$. Inside each KL the $\nu(x_{1:n}) \ln \nu(x_{1:n})$ pieces
cancel, leaving
\begin{align*}
\sum_\nu w_\nu \KL(\nu \,\|\, \rho) - \sum_\nu w_\nu \KL(\nu \,\|\, \xi)
&= \sum_\nu w_\nu \sum_{x_{1:n}}
   \nu(x_{1:n})\,\Bigl[\ln\tfrac{\nu(x_{1:n})}{\rho(x_{1:n})}
                  - \ln\tfrac{\nu(x_{1:n})}{\xi(x_{1:n})}\Bigr] \\
&= \sum_\nu w_\nu \sum_{x_{1:n}}
   \nu(x_{1:n})\,\ln \tfrac{\xi(x_{1:n})}{\rho(x_{1:n})}.
\end{align*}
The factor $\ln(\xi/\rho)$ does not depend on $\nu$, so we swap the order of
summation and pull it out, then use $\sum_\nu w_\nu \nu(x_{1:n}) = \xi(x_{1:n})$:
\begin{align*}
\sum_\nu w_\nu \sum_{x_{1:n}}
   \nu(x_{1:n})\,\ln \tfrac{\xi(x_{1:n})}{\rho(x_{1:n})}
&= \sum_{x_{1:n}} \ln \tfrac{\xi(x_{1:n})}{\rho(x_{1:n})}
   \underbrace{\sum_\nu w_\nu \nu(x_{1:n})}_{=\, \xi(x_{1:n})} \\
&= \sum_{x_{1:n}} \xi(x_{1:n}) \,\ln \tfrac{\xi(x_{1:n})}{\rho(x_{1:n})}
~=~ \KL\bigl(\xi(X_{1:n}) \,\big\|\, \rho(X_{1:n})\bigr).
\end{align*}
Rearranging gives the claimed identity. \qedhere
\end{solution}

\bigskip

\noindent
\textbf{[05]} (Pareto optimality of $\xi$ under KL.) Show that $\xi$ is
Pareto-optimal under $\FKL^n$ in the sense of \cref{def:pareto}: if $\rho$ is
any predictor with
\[
  \FKL^n(\nu, \rho) ~\leq~ \FKL^n(\nu, \xi) \qquad \text{for every } \nu \in \Msol,
\]
then $\rho(x_{1:n}) = \xi(x_{1:n})$ for \emph{every} $x_{1:n} \in \mathbb{B}^n$
(and hence, by chain rule, $\rho(x_t \mid x_{<t}) = \xi(x_t \mid x_{<t})$ at
every history $x_{<t}$ with $\xi(x_{<t}) > 0$).

\textit{Hint.} Multiply the assumed inequality by $w_\nu \geq 0$ and sum over
$\nu \in \Msol$; then apply the KL Pythagorean identity to recognize the extra
non-negative term --- which must vanish.

\begin{solution}
Multiply the assumed inequality by $w_\nu \geq 0$ and sum over $\nu \in \Msol$:
\begin{equation}\label{eq:kl_weighted_dom}
  \sum_\nu w_\nu\,\FKL^n(\nu, \rho)
  ~\leq~
  \sum_\nu w_\nu\,\FKL^n(\nu, \xi).
\end{equation}
The KL Pythagorean identity rewrites the LHS as
\[
  \KL\bigl(\xi(X_{1:n}) \,\big\|\, \rho(X_{1:n})\bigr) \;+\;
  \sum_\nu w_\nu\,\FKL^n(\nu, \xi),
\]
so \cref{eq:kl_weighted_dom} forces
\[
  \KL\bigl(\xi(X_{1:n}) \,\big\|\, \rho(X_{1:n})\bigr) ~\leq~ 0.
\]
KL is non-negative, hence $\KL(\xi(X_{1:n}) \,\|\, \rho(X_{1:n})) = 0$, which
forces $\rho(x_{1:n}) = \xi(x_{1:n})$ for every $x_{1:n} \in \mathbb{B}^n$. By
chain rule, the conditionals agree at every history $x_{<t}$ with
$\xi(x_{<t}) > 0$. \qedhere
\end{solution}

%===================================
\section{Pareto optimality under squared loss}\label{prob:sq}

Squared loss lives one level down: at the conditional distributions
$\nu(\cdot \mid x_{<t})$. The natural Pythagorean decomposition at a fixed
history uses \emph{posterior} weights $w(\nu \mid x_{<t})$ --- those are the
weights that make $\xi(x_t \mid x_{<t})$ the mean of the $\nu(x_t \mid x_{<t})$'s
(via the posterior-predictive form from the recap). The Pareto-optimality
aggregation, however, uses prior weights $w_\nu$. Bridging the two takes one
extra Bayes-rule step.

\bigskip

\noindent
\textbf{[10]} (Per-history Pythagorean identity for squared loss.) Fix any $t
\in \{1, \dots, n\}$ and any history $x_{<t} \in \mathbb{B}^{t-1}$ ---
\emph{the history is arbitrary, not sampled from any environment}. For any
$x_t \in \mathbb{B}$ and any predictor $\rho$, show
\begin{align*}
  &\bigl(\xi(x_t \mid x_{<t}) - \rho(x_t \mid x_{<t})\bigr)^{2}
  ~+~
  \sum_{\nu \in \Msol} w(\nu \mid x_{<t})\,\bigl(\nu(x_t \mid x_{<t}) - \xi(x_t \mid x_{<t})\bigr)^{2} \\
  &=~
  \sum_{\nu \in \Msol} w(\nu \mid x_{<t})\,\bigl(\nu(x_t \mid x_{<t}) - \rho(x_t \mid x_{<t})\bigr)^{2}.
\end{align*}

\textit{Hint.} Add and subtract $\xi(x_t \mid x_{<t})$ inside the squared term,
then expand. Use $\sum_\nu w(\nu \mid x_{<t}) = 1$ and the posterior-predictive
form of $\xi$ from the recap, $\sum_\nu w(\nu \mid x_{<t})\,\nu(x_t \mid x_{<t})
= \xi(x_t \mid x_{<t})$, to kill the cross-term. The weighting must be
\emph{posterior} weights $w(\nu \mid x_{<t})$ (not prior weights $w_\nu$),
because $\xi(x_t \mid x_{<t})$ is the posterior-weighted mean of the
$\nu(x_t \mid x_{<t})$'s --- not the prior one.

\begin{solution}
Throughout, the history $x_{<t}$ and symbol $x_t$ are fixed but arbitrary.
Abbreviate $\nu := \nu(x_t \mid x_{<t})$, $\xi := \xi(x_t \mid x_{<t})$,
$\rho := \rho(x_t \mid x_{<t})$, and $\tilde w_\nu := w(\nu \mid x_{<t})$ for
the duration of the calculation.

\medskip
\textit{Step 1: $\xi$ is the posterior-weighted mean of the $\nu$'s.}
By the posterior-predictive form of the mixture (recap) and the posterior
normalization $\sum_\nu \tilde w_\nu = 1$,
\[
  \sum_\nu \tilde w_\nu \, \nu ~=~ \xi.
\]

\medskip
\textit{Step 2: the cross-term vanishes.}
A direct consequence of Step 1, again using $\sum_\nu \tilde w_\nu = 1$:
\[
  \sum_\nu \tilde w_\nu \, (\nu - \xi)
  ~=~ \sum_\nu \tilde w_\nu \, \nu \;-\; \xi \sum_\nu \tilde w_\nu
  ~=~ \xi - \xi ~=~ 0.
\]

\medskip
\textit{Step 3: expand the square.}
Write $\nu - \rho = (\nu - \xi) + (\xi - \rho)$ and expand:
\begin{align*}
\sum_\nu \tilde w_\nu (\nu - \rho)^2
  &= \sum_\nu \tilde w_\nu\bigl[(\nu - \xi)^2 + 2(\nu - \xi)(\xi - \rho) + (\xi - \rho)^2\bigr] \\
  &= \sum_\nu \tilde w_\nu (\nu - \xi)^2
     \;+\; 2(\xi - \rho)\sum_\nu \tilde w_\nu (\nu - \xi)
     \;+\; (\xi - \rho)^2 \sum_\nu \tilde w_\nu.
\end{align*}
The middle sum is zero by Step 2; the trailing $\sum_\nu \tilde w_\nu = 1$. So
\[
  (\xi - \rho)^2 \;+\; \sum_\nu \tilde w_\nu (\nu - \xi)^2
  ~=~ \sum_\nu \tilde w_\nu (\nu - \rho)^2. \qedhere
\]
\end{solution}

\bigskip

\noindent
\textbf{[05]} (Pareto optimality of $\xi$ under squared loss.) Show that $\xi$
is Pareto-optimal under $\FSQ^n$ in the sense of \cref{def:pareto}: if $\rho$
is any predictor with
\[
  \FSQ^n(\nu, \rho) ~\leq~ \FSQ^n(\nu, \xi) \qquad \text{for every } \nu \in \Msol,
\]
then $\rho(x_t \mid x_{<t}) = \xi(x_t \mid x_{<t})$ at every history $x_{<t} \in
\mathbb{B}^{t-1}$ \emph{with $\xi(x_{<t}) > 0$} (equivalently, $\xi$-almost
surely) and every $x_t \in \mathbb{B}$. Histories that the mixture never reaches
($\xi(x_{<t}) = 0$, i.e., reached by no $\nu \in \Msol$ either) are invisible
to the loss and are not pinned down.

\textit{Hint.} Multiply the assumed inequality by $w_\nu \geq 0$ and sum over
$\nu$; then for each fixed history $x_{<t}$, multiply the per-history
Pythagorean identity from above by $\xi(x_{<t})$ and use the Bayes identity
$w(\nu \mid x_{<t})\,\xi(x_{<t}) = w_\nu \,\nu(x_{<t})$ to bridge from posterior
weights (inside the identity) to prior weights (in the aggregated Pareto
inequality).

\begin{solution}
Multiply the assumed inequality by $w_\nu \geq 0$ and sum over $\nu \in \Msol$:
\begin{equation}\label{eq:sq_weighted_dom}
  \sum_\nu w_\nu\,\FSQ^n(\nu, \rho)
  ~\leq~
  \sum_\nu w_\nu\,\FSQ^n(\nu, \xi).
\end{equation}

Now bridge to the per-history Pythagorean. At each \emph{fixed} history $x_{<t}
\in \mathbb{B}^{t-1}$ and symbol $x_t \in \mathbb{B}$, the identity says
\[
  (\xi - \rho)^2 \;+\; \sum_\nu w(\nu \mid x_{<t})\,(\nu - \xi)^2
  ~=~ \sum_\nu w(\nu \mid x_{<t})\,(\nu - \rho)^2,
\]
with the abbreviations $\nu := \nu(x_t \mid x_{<t})$, $\xi := \xi(x_t \mid x_{<t})$,
$\rho := \rho(x_t \mid x_{<t})$. Multiply by $\xi(x_{<t})$ and use Bayes,
$w(\nu \mid x_{<t})\,\xi(x_{<t}) = w_\nu \,\nu(x_{<t})$:
\[
  \xi(x_{<t})(\xi - \rho)^2
  \;+\; \sum_\nu w_\nu \,\nu(x_{<t})\,(\nu - \xi)^2
  ~=~ \sum_\nu w_\nu \,\nu(x_{<t})\,(\nu - \rho)^2.
\]
Sum over $t = 1, \dots, n$ and all histories $x_{<t} \in \mathbb{B}^{t-1}$ and
symbols $x_t \in \mathbb{B}$. The two $\sum_\nu w_\nu \,\nu(x_{<t})(\cdots)$
terms become exactly $\sum_\nu w_\nu \,\FSQ^n(\nu, \xi)$ and $\sum_\nu w_\nu \,
\FSQ^n(\nu, \rho)$ respectively, so
\[
  \|\xi - \rho\|_\xi^2
  \;+\; \sum_\nu w_\nu\,\FSQ^n(\nu, \xi)
  ~=~
  \sum_\nu w_\nu\,\FSQ^n(\nu, \rho),
\]
where
\[
  \|\xi - \rho\|_\xi^2 ~:=~
  \sum_{t=1}^n \sum_{x_{<t}} \xi(x_{<t}) \sum_{x_t}
    \bigl(\xi(x_t \mid x_{<t}) - \rho(x_t \mid x_{<t})\bigr)^2 ~\geq~ 0.
\]
Combining with \cref{eq:sq_weighted_dom},
\[
  \|\xi - \rho\|_\xi^2 \;+\; \sum_\nu w_\nu\,\FSQ^n(\nu, \xi)
  ~\leq~ \sum_\nu w_\nu\,\FSQ^n(\nu, \xi)
  \;\;\Longrightarrow\;\;
  \|\xi - \rho\|_\xi^2 ~\leq~ 0.
\]
Since $\|\xi - \rho\|_\xi^2$ is a sum of non-negative terms and is itself
$\leq 0$, every term vanishes:
$\xi(x_{<t})\bigl(\xi(x_t \mid x_{<t}) - \rho(x_t \mid x_{<t})\bigr)^2 = 0$
at every $(t, x_{<t}, x_t)$. So $\rho(x_t \mid x_{<t}) = \xi(x_t \mid x_{<t})$
at every history $x_{<t}$ with $\xi(x_{<t}) > 0$. \qedhere
\end{solution}

%===================================
\section*{Closing remarks}

\begin{remark}[Why these are ``Pythagorean'']
Both identities have the shape
\[
  \text{(loss from $\rho$, averaged over $\nu$'s)}
  \;=\;
  \text{(loss from $\xi$, averaged over $\nu$'s)} \;+\;
  \text{(loss from $\rho$ to $\xi$).}
\]
For squared loss this is exactly Pythagoras' theorem in Euclidean space at
each fixed history $x_{<t}$, with $\xi(\cdot \mid x_{<t})$ playing the role
of the orthogonal projection of the ``cloud of $\nu(\cdot \mid x_{<t})$'s''
onto the line of possible reference predictors. For KL it is the Bregman-divergence
analogue at the joint level $x_{1:n}$: $\KL$ behaves like a squared distance,
with $\xi(X_{1:n})$ as the Bregman projection. The crucial property both
proofs share is the ``mean-zero cross-term'': $\sum_\nu w(\nu \mid x_{<t})
(\nu - \xi) = 0$ at the conditional level (squared loss) and $\sum_\nu w_\nu
(\nu(x_{1:n}) - \xi(x_{1:n})) = 0$ at the joint level (KL). In each case the
weighting is whichever one makes $\xi$ the mean of the $\nu$'s --- posterior
weights for the conditional, prior weights for the joint.
\end{remark}

\begin{remark}[Why KL is easier than squared loss]
KL's Pythagorean identity uses prior weights $w_\nu$, matching the weights in
the Pareto sum directly: one identity, one inequality, done. The squared-loss
Pythagorean uses posterior weights $w(\nu \mid x_{<t})$, so bridging to the
prior-weighted Pareto sum costs an extra step (multiply by $\xi(x_{<t})$, apply
Bayes). Geometrically, KL lives at the joint $x_{1:n}$ level where one
identity does everything; squared loss lives at the conditional $(x_t \mid
x_{<t})$ level where the natural mean-zero weighting shifts with the history.
The conclusion is also slightly weaker for squared loss --- $\rho = \xi$ only
at $\xi$-positive histories (i.e., $\xi$-almost surely), not at every history
--- because off-support histories are invisible to a $\nu$-expected loss.
\end{remark}

\begin{remark}[$\xi$ minimizes the prior-weighted loss]
Both identities show more than Pareto optimality. For KL, reading the
joint-level identity as a function of $\rho$, the $\rho$-dependence on the
RHS is a single non-negative term $\KL(\xi(X_{1:n}) \,\|\, \rho(X_{1:n}))$
that vanishes only at $\rho(x_{1:n}) = \xi(x_{1:n})$. So the unique minimizer
of $\rho \mapsto \sum_\nu w_\nu \,\FKL^n(\nu, \rho)$ is $\rho = \xi$. For
squared loss, after the multiply-by-$\xi(x_{<t})$-and-sum aggregation, the
$\rho$-dependence is $\sum_{t,x_{<t},x_t} \xi(x_{<t})(\xi - \rho)^2$, which
vanishes iff $\rho = \xi$ at every $\xi$-positive history. In both cases the
Pareto-optimality result is an immediate corollary.
\end{remark}

\begin{remark}[Versus the standard ``Bayes is admissible'' result]
Pareto optimality is a kind of admissibility statement: nothing weakly
dominates $\xi$ across the class. The standard textbook way to prove this for
Bayesian estimators goes via the complete-class theorem in decision theory and
is considerably more abstract. The Pythagorean identities used here give the
result by direct calculation, with no decision-theoretic machinery and no
assumptions beyond a proper prior ($\sum_\nu w_\nu = 1$). The reason it works
so cleanly is that both losses are convex Bregman divergences in their second
argument --- exactly the setting in which the Bayesian mixture is the unique
minimizer.
\end{remark}

\end{document}
